\chapter{Database deep dive: transactions, migrations, and concurrency}\label{ch:db-deep-dive}

\noindent\textbf{Learning path: Fast path: core.} Learn schema/query invariants needed by persistent applications.\par\medskip

\rustbridgeblock{Function signatures, \code{Result}, \code{Vec<T>}, typed parameters, and \code{?} remain familiar.}{\code{#[query]} plus compiler-owned SQL, typed bind parameters, static row bounds, and \code{Db}/transaction capabilities make database effects reviewable.}{Do not learn string-built SQL as a Velran technique. The new concept is the verified database boundary, not a new way to concatenate queries.}


\invariantblock{Security invariant: data is not SQL structure}{SQL structure remains compiler-owned. User-controlled values enter queries only through typed bind parameters, and result size remains explicitly bounded; validation never licenses string-built SQL.}

\diagnosticblock{Compiler diagnostic}{SEC-A05-001}{The source attempts unsafe SQL construction instead of compiler-owned SQL with typed binds.}{Keep SQL structure static and move request-controlled values into typed bind parameters.}

The previous chapter introduced typed SQL and CRUD. This chapter focuses on the database rules that preserve integrity across concurrent requests and releases. On a first reading, prioritize query contracts, typed binds, transactions, optimistic locking, and migrations; treat capacity, backup, and engine-selection details as later reference.

\section{Three separate responsibilities}

In a Velran application it is useful to separate three layers:

\begin{description}
\item[application code] declares static typed SQL queries, supplies binds, and requests transactions;
\item[runtime] pools connections, binds parameters, applies timeouts and result limits, then validates returned rows against the declared type;
\item[operator] owns database credentials, TLS, database capacity, backups, and the migration process.
\end{description}

The application therefore does not open arbitrary sockets to the database or build its own connection pool. \code{Db} and \code{Transaction} are the capabilities through which database work occurs.

The current V1 data layer also has broad defense-in-depth limits, including a maximum pool size of 32 connections, a 5-second connection-acquire timeout, a 10-second query timeout, a 10,000-row backend result ceiling, and a 16 MiB result-size ceiling. Velran model-list queries are stricter: the compiler requires a literal \code{LIMIT}, and the compiler/runtime cap that list at 1,000 rows. These are safety ceilings, not pagination design; correct queries should still be short, paginated, and indexable.

With multiple server instances, the operator must budget total database connections across all instances:

\begin{lstlisting}
total DB connections ~= number of Velran instances
                       * pool per instance
                       + admin/maintenance reserve
\end{lstlisting}

\section{The query return type is a contract}

Query cardinality is visible in the function type:

\begin{lstlisting}[language=Velran]
#[query]
fn articleById(db: Db, id: i64)
    -> Result<Article, DbError> sql {
    SELECT id, slug, title, body, version
    FROM articles
    WHERE id = :id
}

#[query]
fn articleBySlug(db: Db, slug: Slug)
    -> Result<Option<Article>, DbError> sql {
    SELECT id, slug, title, body, version
    FROM articles
    WHERE slug = :slug
}

#[query]
fn recentArticles(db: Db, offset: i64)
    -> Result<Vec<Article>, DbError> sql {
    SELECT id, slug, title, body, version
    FROM articles
    ORDER BY id DESC
    LIMIT 100 OFFSET :offset
}
\end{lstlisting}

The three forms mean:

\begin{center}
\begin{tabularx}{\textwidth}{@{}L{0.23\textwidth}YY@{}}
\toprule
Type & Expected result & Cardinality violation \\
\midrule
\code{Article} & exactly one row & 0 rows: \code{404}; multiple rows: DB error \\
\code{Option<Article>} & zero or one row & multiple rows: DB error \\
\code{Vec<Article>} & zero or more rows & runtime result limits apply \\
\bottomrule
\end{tabularx}
\end{center}

This is more than convenient notation: result shape is part of the program contract and is enforced at runtime. \code{Option<Article>} does not mean ``take the first row if one exists.'' If the query returns two or more rows, the invariant is broken and the runtime fails closed with a DB error. Similarly, an \code{Article} query returns \code{404 Not Found} for zero rows and does not arbitrarily choose a row if there are several.

A \code{Vec<Model>} query must also declare a static literal row bound in SQL. The bound belongs to the reusable query contract, not to a caller-supplied \code{limit} parameter; use a fixed \code{LIMIT} and paginate with a bounded offset or another explicit access pattern.

\subsection*{Make row shape explicit}

The names and order of \code{SELECT} fields should match the model. Do not use \code{SELECT *} in application queries.

\begin{lstlisting}[language=Velran]
model Article {
    id: i64
    slug: Slug
    title: String
    body: String
    version: i64
}

// good: the same five fields, in the same order
SELECT id, slug, title, body, version
FROM articles
\end{lstlisting}

Explicit projection matters for two reasons. First, the compiler can verify the relationship between model and query. Second, later schema expansion does not silently alter an old query's result shape.

\section{Bind values; do not build SQL}

Client-derived values enter SQL as typed binds:

\begin{lstlisting}[language=Velran]
#[query]
fn findByEmail(db: Db, email: String)
    -> Result<Option<User>, DbError> sql {
    SELECT id, email, display_name
    FROM users
    WHERE email = :email
}
\end{lstlisting}

\code{:email} is not string substitution. The runtime translates it to a backend-specific placeholder and binds it as a parameter.

Do not think in terms of:

\begin{lstlisting}
// wrong approach:
// "SELECT ... ORDER BY " + userOrder
// "WHERE " + userFilter
\end{lstlisting}

Column names, \code{ORDER BY} direction, SQL operators, and entire WHERE fragments must not come from user input. If several sort modes are supported, choose among a small number of static queries in application code.

\section{Read and mutation authority boundaries}

A read query receives a \code{Db} capability. A mutating query receives a \code{Transaction}:

\begin{lstlisting}[language=Velran]
#[query]
fn renameArticle(
    tx: Transaction,
    id: i64,
    title: String
) -> Result<(), DbError> mutates Article by id sql {
    UPDATE articles
    SET title = :title
    WHERE id = :id
}
\end{lstlisting}

The caller opens an explicit transaction block:

\begin{lstlisting}[language=Velran]
#[action]
fn rename(
    ctx: ActionContext,
    db: Db,
    id: i64,
    title: String
) -> Result<Redirect, PageError> {
    let article = articleById(db, id)?;
    authorize article authenticated;
    transaction db {
        renameArticle(tx, article.id, title)?;
    }
    return Ok(redirect(adminArticles()));
}
\end{lstlisting}

For ordinary best-effort application code, a successful block commits and a known database failure rolls back. For retry-sensitive or critical work, do not collapse the outcome to a Boolean success/failure. Capture the transaction outcome explicitly:

\begin{lstlisting}[language=Velran]
let article = articleById(db, id)?;
authorize article authenticated;
let outcome = transaction db {
    renameArticle(tx, article.id, title)?;
};

match outcome {
    Committed => { return Ok(redirect(adminArticles())); }
    RolledBack => { fail conflict; }
    CommitUnknown => { fail conflict; }
}
\end{lstlisting}

\code{CommitUnknown} means the client/runtime did not receive enough evidence to prove whether the commit took effect. It must not automatically trigger a retry. Critical idempotent transactions are compiler-required to capture and exhaustively match this outcome. \code{Transaction} still cannot escape the block, so a closed transaction handle cannot accidentally be reused.

\subsection*{One business operation should be one transaction}

Do not split related mutations across separate actions or separate transactions:

\begin{lstlisting}[language=Velran]
transaction db {
    Article.publishChanged(tx, article.id, version)?;
    audit Article article.id action publish;
        from article.status to ArticleStatus.Published
}
\end{lstlisting}

Here publishing and the business audit trail commit together. If the audit record cannot be written, publishing rolls back too. This prevents a state where the business change happened but its required audit record is missing.

\section{Optimistic locking: do not overwrite someone else's work}

In a CMS or wiki, two editors may open the same record. A simple

\begin{lstlisting}
UPDATE articles SET title = :title WHERE id = :id
\end{lstlisting}

makes the last writer win and can silently discard the first editor's changes.

Use a version column and \code{Changed}:

\begin{lstlisting}[language=Velran]
model Article {
    id: i64
    title: String
    version: i64
}

#[query]
fn updateArticle(
    tx: Transaction,
    id: i64,
    title: String,
    version: i64
) -> Result<Changed, DbError> mutates Article by id sql {
    UPDATE articles
    SET title = :title,
        version = version + 1
    WHERE id = :id AND version = :version
}
\end{lstlisting}

At the call site, load and authorize the article first and pass \code{article.id}, not the request ID directly, into this mutation. The edit form sends the \code{version} value the user saw when opening the page. The \code{Changed} contract is strict: the mutation must affect \textbf{exactly one row}. One changed row is success; zero changed rows become \code{409 Conflict}; two or more rows become a DB error because the query violated its own invariant. \code{Changed} is only valid as the return type of a mutating query, and in V1 it cannot be combined with a \code{RETURNING} clause.

With optimistic locking, if another editor saved in the meantime, \code{WHERE id = :id AND version = :version} matches zero rows and the request becomes a conflict. The correct UX is not an automatic retry. Tell the user that the data changed, load the new version, and let the user decide how to merge.

Typical uses include:

\begin{itemize}
\item wiki pages and CMS articles;
\item tickets or customer records;
\item business objects with workflow state;
\item rarely changed administrator configuration.
\end{itemize}

It is generally unnecessary for append-only events or database-side atomic counters.

\section{Design indexes for access patterns}

An index is useful not because there are many of them, but because it serves a real query.

Natural CMS access patterns might justify:

\begin{lstlisting}[language=SQL]
CREATE UNIQUE INDEX articles_slug_uq
    ON articles(slug);

CREATE INDEX articles_status_id_idx
    ON articles(status, id);
\end{lstlisting}

If the public site loads by slug, make slug unique/indexed. If the admin list filters by status and pages by identifier, a matching composite index may be appropriate.

Some rules:

\begin{itemize}
\item paginate every large list;
\item index frequently used lookup fields;
\item align index order with WHERE/ORDER BY patterns;
\item do not index every column automatically: indexes cost writes and storage;
\item investigate performance with the real backend's \code{EXPLAIN} plan.
\end{itemize}

Velran typed SQL protects query shape; it does not replace the database optimizer.

\section{Migration: schema is part of the release}

In production, the server does not run migrations automatically at startup. This is deliberate: a schema change is a separate operator action with separate credentials and a separate failure boundary.

A typical directory is:

\begin{lstlisting}
migrations/
  0001_create_articles.sql
  0002_articles_slug_index.sql
  0003_articles_add_version.sql
\end{lstlisting}

File names use:

\begin{lstlisting}
NNNN_name.sql
\end{lstlisting}

Do not modify an already-applied migration. The CLI records a SHA-256 checksum and stops on mismatch.

Production commands using the installed CLI:

\begin{lstlisting}
/usr/local/bin/velran-cli migrate status \
  --dir /srv/myapp/current/migrations \
  --db-url-file /run/secrets/migration-db-url

/usr/local/bin/velran-cli migrate verify \
  --dir /srv/myapp/current/migrations \
  --db-url-file /run/secrets/migration-db-url

/usr/local/bin/velran-cli migrate apply \
  --dir /srv/myapp/current/migrations \
  --db-url-file /run/secrets/migration-db-url \
  --lock-timeout-secs 30
\end{lstlisting}

Migration credentials may have more privileges than application DML credentials. The server process should therefore not receive them.

\subsection*{What does the migration CLI guarantee?}

\begin{itemize}
\item applied file name and checksum are verified;
\item an old version number cannot silently receive a new migration later;
\item a database-level migration lock is acquired;
\item migration state advances only after successful application.
\end{itemize}

Backend-specific differences remain. PostgreSQL uses a transaction per migration; SQLite keeps the batch in a transaction using a \code{BEGIN IMMEDIATE} writer lock; MariaDB DDL is not transactional in every case, so partial DDL failure may require operator repair.

\section{Expand/contract: preserve rollback options}

A fast application rollback is straightforward only when the new schema remains compatible with the old application.

For larger changes, use expand/contract:

\begin{enumerate}
\item \textbf{expand}: add a new nullable/defaulted field or a new table while keeping the old application functional;
\item deploy the new application;
\item perform a controlled backfill if necessary;
\item in a later release, \textbf{contract}: remove the old field or constraint after the rollback window has closed.
\end{enumerate}

Immediate rename/drop, a new mandatory field, or an incompatible type change can create a forward-only release. Treat that as an explicit release-review decision.

\section{Handling database errors}

\code{DbError} is a technical error. Do not expose driver messages, SQL text, host names, or connection strings to the end user.

At application level distinguish at least:

\begin{description}
\item[record not found] normal domain state, often handled with an optional query;
\item[409 concurrent modification] optimistic-lock conflict; the user resolves it by reloading/comparing;
\item[database unavailable] infrastructure failure; readiness may report it and the server log can be searched by request ID;
\item[constraint/domain failure] preferably prevent it with validation, while keeping the database constraint as the final integrity boundary.
\end{description}

Do not wrap an entire mutating action in an unlimited automatic retry loop. Repeating a request may repeat side effects, and a transaction conflict is not necessarily transient.

\section{SQLite, PostgreSQL, or MariaDB?}

Velran V1 supports all three backends, but that does not make every SQL dialect feature portable.

\begin{description}
\item[SQLite] excellent for simple or single-host deployments; back up a live application database using an SQLite-aware method rather than plain copying the live file.
\item[PostgreSQL] a strong default for larger or multi-instance production systems; use TLS for remote connections.
\item[MariaDB] also supported, but account for DDL transactional differences especially during migrations.
\end{description}

Portable applications should use a common SQL subset. Use backend-specific features only deliberately, and document the backend dependency in the project.

\section{Backup from a database developer's perspective}

As a database developer, remember one key point here: migration and backup are separate lifecycles. Migration describes how schema changes; backup/restore proves that persistent state can be recovered. For PostgreSQL, MariaDB, and SQLite, use the backend's own consistency-aware backup method and do not assume that a plain file copy is automatically a valid backup.

Full application state -- database, AppFs, local-auth store, recovery manifest, recovery point objective (RPO), recovery time objective (RTO), and restore drills -- is covered in Chapter~\ref{ch:recovery}. It matters here because schema-change planning must already know what state is available for rollback or restoration.

\section{A sample CMS database layer}

A small CMS might start with:

\begin{lstlisting}
Article
  id
  slug          UNIQUE
  title
  body
  status
  hero_image
  version
  created_at
  updated_at

migrations/
  0001_create_articles.sql
  0002_article_indexes.sql
  0003_add_version.sql
\end{lstlisting}

Application patterns:

\begin{itemize}
\item public read: \code{Option<Article>} by slug;
\item admin list: paginated \code{Vec<Article>};
\item edit: \code{Changed} + \code{version};
\item publish/archive: business mutation + audit in one transaction;
\item schema change: separate CLI migration;
\item before deploy: migration verify + backup/restore readiness;
\item after deploy: readiness, logs, metrics, and smoke test.
\end{itemize}

With this model the database layer is not a separate collection of SQL snippets; it becomes part of the web application's type, concurrency, release, and recovery contract.

\section{Chapter checklist}

Before considering a database-backed feature complete, verify:

\begin{itemize}
\item every user value enters SQL as a typed bind;
\item there is no \code{SELECT *}; projections match models;
\item large lists are paginated;
\item important lookup/sort patterns have appropriate indexes;
\item every mutation runs with a \code{Transaction} capability;
\item related mutations share one transaction;
\item human-concurrently-edited records use optimistic locking;
\item already-applied migrations are immutable;
\item migration credentials are not available to the application server;
\item destructive schema changes have a known rollback consequence;
\item DB errors do not leak SQL or credentials to the client;
\item backups have regularly proven restores.
\end{itemize}

\section{Practice: review a query as an interface}
\paragraph{Exercise mode: supported.} Use the supported method defined in the preface.

For each query in a feature, write down its input types, projected columns, cardinality expectation, ordering, and index assumption. Treat that list as an interface contract rather than as an implementation detail hidden inside SQL.

\section{Common mistake: using the database only as storage}
A production database also enforces invariants and concurrency semantics. If correctness depends on a condition remaining true under concurrent requests, encode the authoritative part in the database or transaction design rather than relying only on a prior application-side check.

\section{Running project checkpoint: Secure Notes}
Start from \path{examples/secure-notes/migrations/0001_init.sql}. Review the title/body bounds, publication-state constraint, timestamps, and owner/id index as database-level invariants. Then add pagination while preserving deterministic ordering and principal scoping; document which guarantees belong to SQL/schema design and which remain handler-level policy.
